\begin{table}[t]
\centering
\small
\caption{Table S1. Cost, tax and turnover accounting on the optimised book at \$1bn}
\label{tab:s1}
\begin{tabular}{lrrrrrrrr}
\toprule
Market & Cells & Months & Gross Sharpe & Net Sharpe & After-tax Sharpe & Cost drag (bp/yr) & Tax drag (bp/yr) & Turnover (mean monthly) \\
\midrule
Germany & 8 & 132 & 0.423 & 0.128 & 0.017 & 17.21 & 6.38 & 0.0135 \\
India & 8 & 131 & 0.534 & 0.205 & 0.201 & 9.96 & 0.11 & 0.0067 \\
Japan & 8 & 132 & 0.437 & 0.057 & -0.049 & 42.31 & 11.49 & 0.0290 \\
\bottomrule
\end{tabular}
\par\vspace{2pt}
\parbox{\textwidth}{\footnotesize Optimised book at \$1bn, stage-06 summaries, means over the 8 prediction cells. Net Sharpe and cost drag are from the zero-tax regime rows; after-tax Sharpe and tax drag from the statutory-regime rows. Germany's and India's summaries are verbatim snapshots of earlier untracked outputs (the German one from an older script version with 24 columns); the weight identity was recovered from the weights index. Turnover is endogenous, so the ordering of tax drag by turnover is descriptive and does not separate rate from wedge. No inference is attached.}
\end{table}
